\section{The science studies in full}
\label{sec:ai4s-full}

The general-code studies (Section~\ref{sec:ceiling}) measured the model tier on short general-purpose functions. A scientist would hand
an auditor three different jobs: check the code that computes a result, check that a reported
result follows from that code, and check the data an analysis starts from. We ran one
preregistered study per job, each with ground truth that no model supplies. Scientific code
(\S\ref{sec:ai4s-code}) is judged by scientists' numeric tests, results (\S\ref{sec:ai4s-results})
by re-execution, and data (\S\ref{sec:ai4s-data}) by the faults we injected. Every item in all
three is audited through the product's own audit call with its shipped rulebook, and the same
flag rule holds throughout: an item is flagged when the model grades any finding BLOCKER.

\subsection{Scientific code}
\label{sec:ai4s-code-full}

\textbf{Substrate.} SciCode \citep{tian2024scicode} decomposes 80 research problems from 16
subfields into 338 steps, each a function with numeric tests written by scientists. Its targets
are released separately from the code, which lets the tests act as a hidden oracle. On the
benchmark's development split, the gold code passes 48 of 50 steps in our environment; the two
steps whose gold fails here are excluded, as are the three the official harness skips.

\textbf{Population.} The generator of the general-code studies wrote every step in SciCode's own protocol,
three samples per step, each sample building on its own earlier steps. A candidate is
\emph{correct} if it passes the step's tests and \emph{defective} if it fails or times out. A step
enters either stratum only if at least one sample passes it, so a failure counts only where this
environment has shown the step can be passed. This yields 81 defective instances from 30 problems
and 374 correct ones, from which 150 (53 problems) were drawn by seed.

\textbf{What differs from the general-code studies.} SciCode has no visible and hidden split, and we did not invent
one. The auditor sees the problem, the step, its background, the function header, the allowed
dependencies and the candidate, and no tests of any kind. The shipped cross-vendor auditor and the
generator's own model each read every instance eight times. We reach the generator's model
through a command-line route rather than the API. The route adds a short fixed context and uses
default sampling, so this study's same-vendor family is not temperature-0, unlike the one in the
general-code studies.

\textbf{Two arms.} After the registered arm's outcomes, a second registered arm changed only the
task text, naming the step's function as the deliverable; the program bytes are identical.

\textbf{Residual.} Defective instances flagged by no reading of either family are rated on a
blind sheet under the general-code residual's question (\S\ref{sec:residual}) adapted to science: do the step's prose and background
determine the value the test expects, including tolerance, units, conventions and numerical
method. The raters are the agent that ran the study and a model that audits nothing here.

\textbf{Results.} In these runs, how the task framed the deliverable changed the flag rate
substantially (Figure~\ref{fig:ai4s}a). With the whole problem as the task, the shipped auditor
flagged 87.7\% [72.7, 98.6] of defective and 68.7\% [56.6, 79.7] of correct instances at eight
readings; some of its objections to correct code were that one step did not implement the whole
problem. With the task reframed around the step, flags on correct code fell to 36.0\% [26.4, 45.6], a
paired change of $-32.7$ points [$-42.8$, $-22.4$]. Recall fell too, to 72.8\% [59.3, 87.3], a
change of $-14.8$ points [$-28.6$, $-2.7$], though the cluster sign-flip test for that recall change
gives $p = 0.0625$. The change is the whole reframing (the step named, the problem description moved last
and marked as context, earlier functions exempted), and the arms ran in sequence. The recall of the
generator's own model was 46.9 points below the shipped auditor's in each arm ([$-63.0$, $-30.7$] and
[$-61.5$, $-34.4$]; a route comparison). It flagged 27.3\% [18.4, 37.3] and 18.0\% [11.5, 25.6] of
correct code in the two arms, respectively.
At one reading in the reframed arm the shipped auditor flags 61.7\% [47.9, 75.9] of defective against
19.8\% [12.9, 27.2] of correct instances: it discriminates, modestly. We read a sample of 20 of its
flags on correct code. They range from real divergences between code and step text, through
contradictions inside the benchmark's own instructions, to tests too weak to separate outputs; which
flags point at defects was not established. An independent expert audit of SciCode's test problems found self-contradictory
specifications, non-reproducible gold answers and over-tight tolerances that reject correct solutions
\citep{hu2026scicodeverified}, so some instances labelled defective here may be correct code the
benchmark rejects. Of the 81 defective instances, 7 and 19 escaped all sixteen readings of
the two families in the two arms. Both raters judged most of them (4 of 7, 13 of 19) undetermined by
the prose, but the registered test that this share exceeds the
general-code residual's 77.2\% is killed in both arms by its own interval rule.

\subsection{Scientific results}
\label{sec:ai4s-results-full}

\textbf{Items.} Each correct SciCode instance whose output on one seeded test case is a real
scalar or a real array of at most ten elements becomes a small results increment: the program, the
test's inputs, a script that runs one on the other and writes a run log, a \texttt{results.json}
citing each value's log line, metadata declaring every input at a revision, and a one-paragraph
report whose numbers carry the product's number$\to$source annotations. Seventy-four of the 150
correct instances qualify, from 36 problems. Each yields one clean item and one faulty item. The
faulty item carries one of seven faults, assigned in rotation by seed. Five are \emph{reporting}
faults, where the artefacts disagree with each other or with the log the code wrote: a $\times
10^{3}$ scale slip, a sign flip, a value from another test case, a 5\% relative error, and a
transposed digit pair in the report alone. Two are \emph{fabrications}, where the scale slip or the
5\% error is written consistently into the log, the results and the report, so that every artefact
agrees with every other and none agrees with the code.

\textbf{Three auditors.} The LLM auditor reads the increment four times, with the project's checks
set to the shipped \texttt{science} profile. The model's prompt includes the deterministic
result, as it does in the product, so on reporting faults the model has seen the deterministic
finding. The deterministic profile's own verdict comes from the same call. Re-execution runs the
program on the inputs and compares the output with the reported values.

\textbf{Gate.} Before any model call, every clean item had to pass the science profile with no
BLOCKER and re-execute to its reported value; all 74 did when first built. A two-reading pilot then
found an instrument defect: no file in the item performed the run the report described, and the
auditor blocked on that, correctly. We added the run script, discarded the pilot readings, and
added a third condition to the gate, that each clean item's run script reproduce its run log byte
for byte. All 74 passed all three conditions with no repair.

\textbf{Results.} The deterministic profile does what it was built for (Figure~\ref{fig:ai4s}b).
It flags all 49 reporting faults, because the log, the results and the report no longer agree and
number$\to$source says where, and none of the 25 fabrications, because in those everything agrees
with everything else. It flags no clean item. The LLM auditor flags 23 of the 25 fabrications at
four readings, 21 of them with a sound finding about the fault (post hoc). Those findings work the
result out from the code and inputs, for example evaluating a KL divergence of two small
distributions to 1 where the report says 1000. The other two rest on calculations that are
themselves wrong. Paired on the faulty items, all 23 discordant items go the LLM's way (exact
McNemar $p = 2.4\times10^{-7}$). A problem-level sign-flip test, run post hoc, gives $3.8\times10^{-6}$.
On the reporting faults its rate is of the combined workflow, since its prompt carried the
profile's findings. Re-execution flags every fault except four whose change to a near-zero output falls within
\texttt{np.allclose}'s default absolute tolerance; two of those outputs are floating-point zeros.

The LLM also flags 12 of the 74 clean items. We read each of these flags (post hoc), and 8 are right. In 9 instances the
benchmark passes the test's inputs inside the call itself, so the report's sentence that the
program ran ``on the inputs in \texttt{inputs.py}'' was false; 7 of those clean items were flagged
for it. Of four outputs whose unit the solution documents, two were flagged for our
\texttt{dimensionless} label; one of the two is also among the seven. The remaining 4 flags rest on wrong calculations or, in one case, on
a syntax error that is not there. On this evidence a provenance chain and a reader that recomputes
are complements: the first is exact about agreement and blind to consistent error, and the second
sees consistent error but is sometimes wrong itself.



\subsection{Scientific data}
\label{sec:ai4s-data-full}

\textbf{Substrate.} We use four physical-science tables from the UCI repository, each under CC BY 4.0:
concrete compressive strength \citep{yeh1998concrete}, airfoil self-noise
\citep{brooks1989airfoil}, combined-cycle power-plant output \citep{tufekci2014ccpp}, and
superconductor critical temperature \citep{hamidieh2018superconductor}. Exact duplicate rows and
physically impossible rows (relative humidity above 100\%) were removed from each pool before any
item was drawn.

\textbf{Items.} Each item is 120 rows sampled from one pool with a data card that gives every
column's meaning and unit and, where the source states one, its range. There are 35 clean items
per dataset and five per dataset for each of seven faults: mixed units in 20\% of one column's
rows, three negated values of a non-negative quantity, 12 duplicated rows, the target shuffled
among 30\% of rows, two columns with disjoint ranges swapped, a $-999$ sentinel in 5\% of one
column's cells, and a continuous column rounded to integers.

\textbf{Auditors.} Both model families read each item four times, with the dataset as the
deliverable and the card as its documentation. The comparison is a deterministic validator written
from the card alone and frozen before any item was built: declared ranges, non-negativity,
duplicates, sentinels and integer-valued continuous columns. By construction it catches the
negations, duplicates, sentinels and rounding. It sees mixed units or swapped columns only when
they push a value below zero or out of a declared range, or make a continuous column
integer-valued, and a shuffled target not at all. The question is whether the model catches
what the validator cannot. Two instrument defects were found and fixed by amendment before the
readings reported here: the card and file named one column differently, and the files first sat
outside the audited scope, so the product had told the model there was nothing to audit.

\textbf{Results.} The two tiers flag overlapping but different faulty items (Figure~\ref{fig:ai4s}c). The validator flags
92 of 140 faulty items (65.7\% [57.9, 73.6]) and no clean item. The shipped auditor flags 65 of 140
(46.4\% [38.6, 54.3]) at 4 of 140 clean items (2.9\% [0.7, 5.7]), and their union flags 112 of 140
(80.0\% [73.6, 86.4]) at the same 4. These intervals resample items within dataset; four datasets
are too few to resample, so no interval here carries between-dataset variation. Of the 20 faulty
items the auditor flags and the validator misses, 18 carry a finding about the injected fault
itself, 9 of them mixed units and 9 swapped columns (post hoc). In one, the auditor reasoned that
three atoms spanning a 135\,pm radius range cannot have a mean radius of 1.6\,pm, which is the cell
the fault had converted to {\aa}ngstr\"oms. The other 2, and all 4 flags on clean items, are requests for
fuller column descriptions on one dataset's card; so are 2 more of its faulty-item flags, which the
validator also caught. Counting only findings about the injected fault (post hoc), the auditor flags
61 of 140 and the union 110, at no clean item. The auditor flagged no shuffled target and no
duplicated rows (0 of 20 each), and it did worse than the validator on every fault the validator
was written to catch. The generator's own model, read through its command-line route, is not a usable
data auditor here: it flags 67 of 140 clean items, and in at least 48 of those one finding states a
row count, field count, missing column or truncation that the file contradicts (post hoc lower
bound).

The deterministic validator is not neutral. It was frozen from the card before any item existed,
but it omits one range the concrete source's documentation states and uses a list of continuous
columns the card does not give; with the omitted range it would flag 93 of 140. Unit conversion
also left floating-point tails that mark some converted cells in two of the tables. None of the
auditor's mixed-unit findings cites them, but that does not show they had no influence.
